%=============================================================================!
% 4.A  ANSWERS TO THE CHECKS                                                  !
%=============================================================================!
\unnumbered{Answers to the checks}
\label{sec:os-answers}

\answerto{sec:os-circle}When $\omega t + \phi = \pi/2$ the point $P$ is at the top
of the reference circle, so the block is at the centre, $x = A\cos(\pi/2)
= 0$. Its velocity is $-A\omega\sin(\pi/2) = -0.05\times10 =
\SI{-0.5}{\metre\per\second}$: it moves at its greatest speed towards
negative $x$, as $P$, at the top and going anticlockwise, moves to the
left.

\answerto{sec:os-energy}Where $U = K$, each is $E/2$, so $\frac12 kx^2 =
\frac12\times\frac12 kA^2$ and $x = \pm A/\sqrt2 = \pm4/\sqrt2 =
\pm\SI{2.83}{\centi\metre}$.

\answerto{sec:os-phase}The amplitude is $A = \sqrt{x_0^2 + (v_0/\omega)^2}
= 0.3/10 = \SI{0.03}{\metre} = \SI{3}{\centi\metre}$ by \cref{sec:ne-spring}, and the largest
speed is $A\omega = \SI{0.3}{\metre\per\second}$, so the semi-axes are
\SI{3}{\centi\metre} and \SI{0.3}{\metre\per\second}. It starts at the top
of the ellipse, $(0, \SI{0.3}{\metre\per\second})$, and moves clockwise.

\answerto{sec:os-small}The period $2\pi\sqrt{L/g}$ is proportional to
$\sqrt L$, so four times the length doubles it. The mass cancelled from
$\omega_0 = \sqrt{g/L}$, so it does not matter.

\answerto{sec:os-anharmonic}Beyond small angles the restoring force,
proportional to $\sin\theta$, grows more slowly than $\theta$, so over most
of a large swing the bob is pulled back less strongly than a spring would
pull it. At the ends of a \ang{90} swing the pull along the arc is
$mg\sin\ang{90} = mg$, against $mg\,\pi/2$ for the spring of stiffness
$mg/L$ stretched by $L\pi/2$, so the bob turns round more slowly and
spends longer there. The extra time spent near the ends outweighs the
higher speed at the bottom; \cref{eq:os-period} gives 1.180 times the
small-swing period, against 1.002 for a \ang{10} swing.

\answerto{sec:os-damped}Critical damping needs $\frac12\gamma = \omega_0$,
so $\gamma = \SI{20}{\per\second}$. With $\gamma = \SI{8}{\per\second}$,
$\omega_{\mathrm d} = \sqrt{\omega_0^2 - \frac14\gamma^2} = \sqrt{100 -
16} = \SI{9.165}{\radian\per\second}$,
and the period is $2\pi/9.165 = \SI{0.686}{\second}$.

\answerto{sec:os-driven}At $\omega = \omega_0$, $A = f/(\gamma\omega_0) =
1/(0.5\times10) = \SI{0.2}{\metre} = \SI{20}{\centi\metre}$. For a very slow push, $A \to
f/\omega_0^2 = 1/100\ \si{\metre} = \SI{1}{\centi\metre}$, the stretch of
the spring under the force held still.

\answerto{sec:os-bodies}$m_{\mathrm r} = m^2/(2m) = m/2$, and the separation
oscillates with $\omega = \sqrt{k/m_{\mathrm r}} = \sqrt{2k/m}$.

\answerto{sec:os-modes}In the slow mode both carts move by the same amount
at every instant, so the distance between them never changes and the
middle spring is never stretched: it exerts no force, and each cart swings
on its outer spring alone, with $\omega = \sqrt{k/m}$.

\answerto{sec:os-molecules}$\nu = 2.99792458\times10^{-5}\times1000 =
\SI{0.02998}{\per\femto\second}$, so the period is $1/\nu =
\SI{33.36}{\femto\second}$.
